% BHU PYQ -- Manually transcribed & error-corrected
% Paper: BCH-114 : Monetary Theory
% Exam: B.Com. (Hons.) Semester I Examination, 2022-23

\documentclass[11pt,a4paper]{article}
\usepackage[a4paper,top=2.5cm,bottom=2.5cm,left=2.8cm,right=2.8cm,headheight=14pt]{geometry}
\usepackage{amsmath,amssymb}
\usepackage[T1]{fontenc}
\usepackage[utf8]{inputenc}
\usepackage{lmodern,microtype,fancyhdr,enumitem,hyperref,lastpage,parskip}

\hypersetup{
    colorlinks=true,
    urlcolor=black,
    linkcolor=black,
    pdftitle={BCH-114: Monetary Theory -- BHU B.Com. (Hons.) Sem I 2022-23}
}

\pagestyle{fancy}
\fancyhf{}
\renewcommand{\headrulewidth}{0.4pt}
\renewcommand{\footrulewidth}{0.4pt}
\fancyhead[L]{\small\textit{Banaras Hindu University}}
\fancyhead[R]{\small\textit{BCH-114: Monetary Theory}}
\fancyfoot[C]{\small Page \thepage\ of \pageref{LastPage}}

\newlist{parts}{enumerate}{2}
\setlist[parts,1]{label=(\alph*),leftmargin=*,itemsep=5pt,topsep=3pt}
\setlist[parts,2]{label=(\roman*),leftmargin=*,itemsep=3pt,topsep=2pt}
\newcommand{\pts}[1]{\hfill\textit{[#1]}}

\begin{document}

\begin{center}
  {\large\bfseries BANARAS HINDU UNIVERSITY}\\[2pt]
  {\normalsize B.Com. (Hons.) --- Semester I Examination, 2022-23}\\[6pt]
  \rule{0.8\linewidth}{0.5pt}\\[6pt]
  {\large\bfseries Paper No. BCH-114: Monetary Theory}\\[4pt]
  {\normalsize Time: 3 Hours\hfill Maximum Marks: 70}
\end{center}

\rule{\linewidth}{0.4pt}

\smallskip
\noindent\textit{Instructions: Answer all questions. All questions carry equal marks (14 marks each).}

\bigskip

\noindent\textbf{Question 1.}
Define money and discuss its main approaches.\pts{14}

\centerline{\textbf{OR}}

Discuss the importance of money in capitalist and socialist economy.\pts{14}

\medskip\rule{\linewidth}{0.2pt}

\noindent\textbf{Question 2.}
Discuss the different kinds of monetary standard.\pts{14}

\centerline{\textbf{OR}}

Explain the salient features, merits and demerits of gold standard.\pts{14}

\medskip\rule{\linewidth}{0.2pt}

\noindent\textbf{Question 3.}
What is bank money? Describe the process and limitations of credit creation by banks.\pts{14}

\centerline{\textbf{OR}}

What is meant by value of money? Discuss the Fisher's quantity theory of money.\pts{14}

\medskip\rule{\linewidth}{0.2pt}

\noindent\textbf{Question 4.}
Explain the Cambridge equation of the quantity theory of money. Is it an improvement over Fisher's equation? Describe.\pts{14}

\centerline{\textbf{OR}}

Explain the various methods for the issue of paper currency notes and which method is adopted by the Reserve Bank of India?\pts{14}

\medskip\rule{\linewidth}{0.2pt}

\noindent\textbf{Question 5.}
Why is inflation caused? What are the various measures usually adopted for controlling inflation? Explain.\pts{14}

\centerline{\textbf{OR}}

Distinguish between disinflation and reflation and explain the causes and effects of both.\pts{14}

\vfill
\rule{\linewidth}{0.4pt}
\begin{center}\small
\href{https://bhu-exam.vercel.app/}{Click here to visit our website}\\[4pt]
\href{https://me-aryan.vercel.app/}{Click here to visit our contributor}\\[4pt]
\href{https://quashan-me.vercel.app/}{Click here to visit our contributor}
\end{center}

\end{document}
